\documentclass[10pt]{article}
\usepackage[utf8]{inputenc}
\usepackage[T1]{fontenc}
\usepackage{amsmath}
\usepackage{amsfonts}
\usepackage{amssymb}
\usepackage[version=4]{mhchem}
\usepackage{stmaryrd}
\usepackage{bbold}

\begin{document}
$L 3$

\section*{L3: Greatest Common divisors}
\begin{itemize}
  \item Given $a, b \in \mathbb{Z}$, both not zero, consider
\end{itemize}

$$
\delta:=\{c: c \in \mathbb{Z}, c \mid a \text {, and } c \mid b\} \text {. }
$$

The set $\delta$ is non-empty since $\pm 1 \in S$, and finite since $O$ is the only integer that has an infinite number of divisors, and at least one of $a$ and $b$ are non-zero. Thus it makes sense to speak of the greatest element of $S_{\text {, }}$ and that element is necessarily positive.\\
Definition: Let $a, b \in \mathbb{Z}$ with $a$ and $b$ not both zero. The greatest common divisor of a and $b$, denoted $(a, b)$, is the greatest positive integer $d$ such that dla and dib. If $(a, b)=1$, then $a$ and $b$ are said to be relatively prime.\\
Remark: Note that $(0,0)$ is undefined.\\
Remark: If $(a, b)=d$, then

$$
(-a, b)=(a,-b)=(-a,-b)=d . \quad \text { (*) }
$$

Thus without loss of generality we vestrict our\\
computations of the gcd of two non-negative integers in the following example.\\
Example:

\begin{itemize}
  \item Divisors of $24: \pm 1, \pm 2, \pm 3, \pm 4, \pm 6, \pm 8, \pm 12, \pm 24$
\end{itemize}

Divisors of $60: \pm 1, \pm 2, \pm 3, \pm 4, \pm 5, \pm 6, \pm 10$,

$$
\pm(2), \pm 15, \pm 20, \pm 30, \pm 60
$$

Thus $(24,60)=12$.

\begin{itemize}
  \item In general, if $a \in \mathbb{Z}$ with $a \neq 0$, then $(a, 0)=|a| . \quad(* *)$
  \item Divisors of $35: \pm 1, \pm 5, \pm 7, \pm 35$.
\end{itemize}

Thus $(24,35)=1$, and 24 and 35 are velatively prime.\\
Remark: If $(a, b)=d$, then $d \mid a$ and $d \mid b$.\\
Proposition 1.10: Let $a, b \in \mathbb{Z}$ with $(a, b)=d$. Then $(a / d, b / d)=1$.\\
Proof: Let $d^{\prime}:=(a / d, b / d)$. Then $d^{\prime} \mid a / d$ and $d^{\prime} \mid b / d$, so there exist $e, f \in \mathbb{Z}$ such\\
that $a / d=d^{\prime} e$ and $b / d=d^{\prime} f$. Thus $a=d^{\prime} d e$ and $b=d^{\prime} d f$, and consequently we have $d^{\prime} d \mid a$ and $d^{\prime} d \mid b$. This implies that $d^{\prime} d$ is a common divisor of $a$ and $b$, and since $d$ is the $\operatorname{gcd}$ of $a$ and $b$, we have $d^{\prime}=1$, as required.

Proposition 1.11: Let $a, b \in \mathbb{Z}$ with $a$ and $b$ not both zero. Then\\
$(a, b)=\min \{m a+n b: m, n \in \mathbb{Z}, m a+n b>0\}$.\\
Proof: We first need to prove that the minimum exists. Observe that

$$
\{m a+n b: m, n \in \mathbb{Z} \partial m a+n b>0\} \neq \varnothing,
$$

Since, without loss of generality, $a \neq 0$, and then either $l a+O b>0$ or $-\mid a+O b>0$. Thus min $\{m a+n b: m, n \in \mathbb{Z}, m a+n b>0\}$ exists by the well-ordering principle, denote it $d$. Then $d=m^{\prime} a+n^{\prime} b$ for some $m^{\prime} n^{\prime} \in \mathbb{Z}$.

We first prove that dla and d/b. By the (4) division algorithm there exist $q, r \in \mathbb{Z}$ such that

$$
a=d q+r, \quad 0 \leqslant r<d
$$

Now,

$$
\begin{aligned}
r=a-d q & =a-\left(m^{\prime} a+n^{\prime} b\right) q \\
& =\left(1-q m^{\prime}\right) a-q n^{\prime} b_{J}
\end{aligned}
$$

and we have that $n$ is an integral linear combination of $a$ and $b$. Since $O \leq r<d$ and $d$ is the minimum positive integral linear combination of $a$ and $b$, we have $r=0$, from which $a=d q+r$ implies that $a=d q$. Thus $d l a$, and by a similar argument, $d \mid b$. It remains to show that $d$ is the greatest common divisor. Let $c$ be any common divisor of $a$ and $b$, then $c l a$ and $c l b$. Then $c\left(m^{\prime} a+n^{\prime} b=d\right.$ by Proposition 1.2, and so $c \leqslant d$.

\begin{itemize}
  \item Proposition 1.11 tells us that $(a, b)=d$ implies that d may be expressed as an integral linear combination of $a$ and $b$. The converse of this statement is true if $d=1$, namely, if 1 is expressible as an integral linear combination of $a$ and $b$, then $(a, b)=1$. Why?\\
Well, $1=m a+n b$ for some $m, n \in \mathbb{Z}$, and any $d \in \mathbb{Z}$ such that $d \mid a$ and $d \mid b$ must then satisfy $d \mid m a+n b=1$ by Proposition $(.2$, so $(a, b)=1$.
  \item Proposition 1.11 also can be used to show that a common divisor is not only less them the greatest common diviser, but also divides the greatest common divisor:\\
$(a, b)=a m+b n$ for some $m, n \in \mathbb{Z}$, so if\\
dla and $d \mid b$, then $d \mid a m+b n=(a, b)$ by Proposition 1.2.
\end{itemize}

Definition: Let $a_{1}, a_{2}, \ldots, a_{n} \in \mathbb{Z}$ with $a_{1}, a_{2}, \ldots, a_{n}$ not all zero. The greatest common divitor of $a_{1}, a_{2}, \ldots, a_{n}$, denoted $\left(a_{1}, a_{2}, \ldots, a_{n}\right)$ is the greatest integer $d$ such that $d \mid a_{i}$ for all $i=1, \ldots, n$. If $\left(a_{1}, a_{2}, \ldots, a_{n}\right)=I$, then $a_{1}, a_{2}, \ldots, a_{n}$, are said to be velatively prime. If $\left(a_{i}, a_{j}\right)=1$ for all pairs $i, j$ with $i \neq j$, then $a_{1}, a_{2}, \ldots, a_{n}$ are said to be pairwise relatively prime.\\
Example: The common dividoss of $24,60,30$ are $\pm 1, \pm 2, \pm 3, \pm 6$. Thus $(24,60,30)=6$.

Pairwise relatively prime integers are relatively prine, but the converse is not always true.\\
Example: Check that $(24,60,49)=1$, so $24,66,49$ are relatively prime. However, $(24,60)=12 \neq 1$, so they are not pairwise

\section*{relatively prime.}
\section*{The Enclidean Algorithm}
\begin{itemize}
  \item We want a move efficient method for computing the god of two integers than just listing of their divisors.\\
Lemma 1.12: If $a, b \in \mathbb{Z}, a \geqslant b>0$, and $a=b q+r$ with $q, r \in \mathbb{Z}$, then $(a, b)=(b, r)$.\\
Proof: Let $c$ be a common divisor of a and $b$. Then $c / a$ and $c / b$ imply that $c \mid a-q b=r$, and so $c$ is a common divisor of $b$ and $r$.
\end{itemize}

Now let $c$ be a common divisor of $b$ and $r$, so that $c \mid b$ and $c \mid r$. Then $c \mid q b+r=a$ by Proposition 1.2 , from which $c(a)$ and so $c$ is a common divisor of $a$ and $b$.\\
The set of common divisors of $a$ and $b$ is the same as the set of common divisors of $b$ and $r$, and so $(a, b)=(b, r)$.

\begin{itemize}
  \item The Enclidean algorithm for computing\\
the ged of two positive integers comes from Book VII of Enclid's "Elements".\\
Theorem 1.13: Let $a, b \in \mathbb{Z}$ with $a \geqslant b>0$. By the division algorithm (Theorem 1.4), there exist $\left.q_{1}\right) r_{1} \in \mathbb{Z}$ such that
\end{itemize}

$$
a=b q_{1}+r_{1}, 0 \leqslant r_{1}<b .
$$

If $r_{1}>0$, there exist $q_{2}, r_{2} \in \mathbb{Z}$ by the division algosithm such that

$$
b=r_{1} q_{2}+r_{2}, \quad 0 \leqslant r_{2}<r_{1} .
$$

If $r_{2}>0$, there exist $q_{3}, r_{3} \in \mathbb{Z}$ such that

$$
r_{1}=r_{2} q_{3}+r_{3}, \quad 0 \leq r_{3}<r_{2}
$$

Continue this process. Then $r_{n}=0$ for some $n \geqslant 1$, and

$$
(a, b)= \begin{cases}b & \text { if } n=1 \\ r_{n-1} & \text { if } n>1\end{cases}
$$

Proof: Note that $r_{1}>r_{2}>r_{3}>\ldots$ If $r_{n} \neq 0$ for all $n \geqslant 1$, then $r_{1}, r_{2}, r_{3}, \ldots$ is an\\
infinite, strictly decreasing sequence of positive integers, which is not possible!\\
Thus $r_{n}=0$ for some $n$. If $n>1$ we repeatedly apply Lemma 1.12 to obtain\\
$(a, b)=\left(b, r_{1}\right)=\left(r_{1}, r_{2}\right)=\left(r_{2}, r_{3}\right)=\cdots=\left(r_{n-1}, r_{n}\right)$

$$
\left.=\left(r_{n-1}, 0\right)=r_{n-1}\right)
$$

as desired. If $n=1$, the statement is clear.

Remark: The Enclidean Algosithm finds the ged of $a, b$ when they are positive integers. By (*) and (**), all other cases are reducible to this case or trivial.\\
Example: Find $(803,154)$ using the Euclidean algorithm:


\begin{align*}
803 & =154 \cdot 5+33  \tag{1}\\
154 & =33 \cdot 4+22  \tag{2}\\
33 & =22 \cdot 1+11  \tag{3}\\
22 & =11 \cdot 2+0 \tag{4}
\end{align*}


Thus $(803,154)=11$.

\begin{itemize}
  \item Recall that Rroposition 1.11 says that the ged of two integen is expressible as an integral linear combination of the two integers. The Euclidean algorithm gives a Systematic way to find Such a linear combination.\\
Example: Express $(803,154)$ as an integral linear combination of 803 and 154.\\
We work through the previous example backward. We have
\end{itemize}


\begin{align*}
(803,154)=11 & =33-22 \text { by } \\
& =33-(154-33 \cdot 4) \text { by }(2) \\
& =33 \cdot 5-154 \\
& =(803-154 \cdot 5) 5-154 \text { by }(0)  \tag{0}\\
& =803 \cdot 5-154 \cdot 26 \\
& =5.803+(-26) \cdot 154 \text {. }
\end{align*}


Remark: This integral linear combination of\\
$(803,154)$ is not unique. Can verify that $11=19.803-99.154$ is another such expression. We will study this more later in the course.

\begin{itemize}
  \item No integral linear combination of 803 and
\end{itemize}

154 can yield any of the integers $1,2,3, \ldots, 10$ by Proposition 1.11.


\end{document}